% Dibangkitkan oleh revisi/hasil/gerbang_batch/bab5_jalan.py. Jangan diubah dengan tangan.
\begin{table}[htbp]
\caption{Korelasi Spearman selisih $J$ ARLS-PPO terhadap Fixed-L3 + penjaga dengan sifat instans uji}
\label{tab:uji-sifat-fix3-g}
\centering\footnotesize
\setlength{\tabcolsep}{4pt}
\begin{adjustbox}{max width=\textwidth}
\begin{tabular}{lrrrr}
\toprule
Sifat instans & \multicolumn{2}{c}{$\lambda=0{,}3$} & \multicolumn{2}{c}{$\lambda=3$} \\
 & uk.~1 & uk.~2 & uk.~1 & uk.~2 \\
\midrule
Banyak tugas & $+0{,}30$ & $-0{,}25$ & $+0{,}33$ & $-0{,}24$ \\
Banyak pengembang & $+0{,}19$ & $+0{,}17$ & $+0{,}37$ & $+0{,}24$ \\
Tugas per pengembang & $+0{,}13$ & $-0{,}42$ & $+0{,}05$ & $-0{,}43$ \\
Fleksibilitas (pengembang per tugas) & $+0{,}57$ & $+0{,}31$ & $+0{,}17$ & $-0{,}05$ \\
Beban per pengembang & $+0{,}13$ & $-0{,}33$ & $+0{,}01$ & $-0{,}44$ \\
Ketimpangan keahlian & $+0{,}31$ & $+0{,}06$ & $+0{,}27$ & $+0{,}00$ \\
Keketatan tenggat & $-0{,}05$ & $-0{,}13$ & $-0{,}11$ & $-0{,}17$ \\
Prasyarat per tugas & $+0{,}28$ & $-0{,}27$ & $+0{,}29$ & $-0{,}21$ \\
\bottomrule
\end{tabular}
\end{adjustbox}
\par\smallskip\footnotesize\textit{Catatan:} Satu titik per instans uji (38 instans). Positif berarti makin besar sifat itu, makin besar selisih $J$ (ARLS-PPO makin tertinggal). Deskriptif, tanpa uji hipotesis.
\end{table}
